%! Vectors — Unit 4, Lesson 4.8 — Guided Notes (Answer Key)
\documentclass[10pt]{article}
\usepackage{precalculus-article}
\usepackage{precalculus-key}   % pulls in -boxes; do NOT also load -boxes
\newcommand{\TallMath}[1]{$\displaystyle #1\rule[-1.4em]{0pt}{3.2em}$}

\begin{document}

\pageheader{Unit 4, Lesson 4.8}{Guided Notes --- Answer Key}
\namedateperiod

\vspace{0.1in}

\begin{objectivebox}
By the end of this lesson, I will be able to:
\begin{itemize}[leftmargin=1.5em, itemsep=3pt]
  \item \ansline{Find the component form and magnitude of a vector from its tail and head.}
  \item \ansline{Perform scalar multiplication and vector addition; compute the dot product.}
  \item \ansline{Find the unit vector in the direction of a given vector.}
  \item \ansline{Use the dot product formula to find the angle between two vectors and test for perpendicularity.}
\end{itemize}
\end{objectivebox}

\vspace{0.08in}

\begin{vocabbox}
\termblanklong{vector} \ans{A directed line segment described by magnitude and direction; written $\langle a,b\rangle$.}
\termblanklong{tail / head} \ans{Starting point (tail) and ending point (head) of the directed segment.}
\termblanklong{magnitude} \ans{$\|\vec{v}\|=\sqrt{a^2+b^2}$; the length of the vector.}
\termblanklong{scalar multiplication} \ans{$k\langle a,b\rangle=\langle ka,kb\rangle$; scales the magnitude by $|k|$.}
\termblanklong{dot product} \ans{$\langle a_1,b_1\rangle\cdot\langle a_2,b_2\rangle=a_1 a_2+b_1 b_2$; a scalar.}
\termblanklong{unit vector} \ans{A vector of magnitude 1; $\hat{u}=\vec{v}/\|\vec{v}\|$.}
\end{vocabbox}

\vspace{0.08in}

\begin{hookbox}
Two drones receive identical displacement instructions. Drone A leaves $(1,2)$ and
arrives at $(5,6)$. Drone B leaves $(3,0)$ and arrives at $(7,4)$.
How far east and how far north does each drone travel? Why do they receive the same instructions?

\ansline{Each drone moves 4 units east and 4 units north; the displacement vector $\langle 4,4\rangle$ is the same regardless of starting position.}
\end{hookbox}

\vspace{0.08in}

\begin{notesbox}{Section 1 --- Vector Components and Magnitude}
\small
\textbf{Definition.} A \textbf{vector} from tail $P_1=(x_1,y_1)$ to head $P_2=(x_2,y_2)$ has
component form $\vec{v} = \langle a, b\rangle$ where:

\vspace{0.06in}
\hspace{1em} $a = x_2 - x_1 = $ \ans{$x_2 - x_1$} \qquad $b = y_2 - y_1 = $ \ans{$y_2 - y_1$}

\vspace{0.08in}
\textbf{Magnitude:} $\|\vec{v}\| = \sqrt{a^2 + b^2} = $ \ans{$\sqrt{(x_2-x_1)^2+(y_2-y_1)^2}$}

\vspace{0.12in}
\textbf{Direction from trig.} If $\|\vec{v}\| = r$ and direction angle is $\theta$:

\vspace{0.06in}
\hspace{1em} $a = r\cos\theta = $ \ans{$r\cos\theta$} \qquad $b = r\sin\theta = $ \ans{$r\sin\theta$}

\vspace{0.12in}
\textbf{Example.} Tail $P_1=(1,3)$, head $P_2=(5,6)$.

\vspace{0.06in}
\hspace{1em} $a = $ \ans{$4$} \quad $b = $ \ans{$3$} \quad
$\vec{v} = \langle$ \ans{$4$}$,$ \ans{$3$}$\rangle$ \quad
$\|\vec{v}\| = $ \ans{$\sqrt{16+9}=5$}
\end{notesbox}

\vspace{0.08in}

\begin{notesbox}{Section 2 --- Scalar Multiplication and Vector Addition}
\small
\textbf{Scalar multiplication.} $k\langle a,b\rangle = \langle$ \ans{$ka,\, kb$}$\rangle$

\vspace{0.06in}
New magnitude: $|k| \cdot \|\vec{v}\|$. \quad
Direction: \ans{same} when $k>0$; \ans{reversed} when $k<0$.

\vspace{0.12in}
\textbf{Vector addition.} $\langle a_1,b_1\rangle + \langle a_2,b_2\rangle
= \langle$ \ans{$a_1+a_2,\; b_1+b_2$}$\rangle$

\vspace{0.06in}
Geometric interpretation: place tail of second vector at \ans{head} of first;
resultant runs from \ans{first tail} to \ans{second head}.

\vspace{0.12in}
\textbf{Example.} $\vec{u}=\langle 3,1\rangle$, $\vec{w}=\langle -1,4\rangle$.

\vspace{0.06in}
\renewcommand{\arraystretch}{1.6}
\begin{tabular}{lll}
$\vec{u}+\vec{w} = \langle$\ans{$2$}$,$\ans{$5$}$\rangle$ &\quad&
$2\vec{u} = \langle$\ans{$6$}$,$\ans{$2$}$\rangle$ \\
\end{tabular}
\end{notesbox}

\vspace{0.08in}

\begin{notesbox}{Section 3 --- Dot Product and Unit Vectors}
\small
\textbf{Dot product.}
$\langle a_1,b_1\rangle \cdot \langle a_2,b_2\rangle
= a_1 a_2 + b_1 b_2 = $ \ans{$a_1 a_2 + b_1 b_2$}
\quad (result is a \ans{scalar}, not a vector)

\vspace{0.12in}
\textbf{Unit vector.} A vector of magnitude \ans{$1$}.

\vspace{0.06in}
$\hat{u} = \dfrac{\vec{v}}{\|\vec{v}\|}
= \left\langle \dfrac{a}{\|\vec{v}\|},\, \dfrac{b}{\|\vec{v}\|} \right\rangle$

\vspace{0.1in}
\textbf{Standard basis vectors:}
$\hat{\imath} = \langle$ \ans{$1,0$}$\rangle$ \quad
$\hat{\jmath} = \langle$ \ans{$0,1$}$\rangle$

Any vector: $\langle a,b\rangle =$ \ans{$a$}$\hat{\imath}\, +\, $\ans{$b$}$\hat{\jmath}$

\vspace{0.12in}
\textbf{Example.} $\vec{v}=\langle 3,4\rangle$, $\|\vec{v}\| = $ \ans{$5$}

\vspace{0.06in}
$\hat{u} = \left\langle \dfrac{3}{\ans{5}},\, \dfrac{4}{\ans{5}} \right\rangle
= \langle$\ans{$3/5$}$,$\ans{$4/5$}$\rangle$ \quad
Verify: $\|\hat{u}\| = $ \ans{$\sqrt{9/25+16/25}=1$}
\end{notesbox}

\vspace{0.08in}

\begin{notesbox}{Section 4 --- Angle Between Vectors \& Perpendicularity}
\small
\textbf{Angle formula.}
$\vec{v}_1 \cdot \vec{v}_2 = \|\vec{v}_1\|\,\|\vec{v}_2\| \cos\theta$

\vspace{0.08in}
Solving for $\theta$:
\quad $\theta = \cos^{-1}\!\left( \dfrac{\vec{v}_1 \cdot \vec{v}_2}{\|\vec{v}_1\|\,\|\vec{v}_2\|} \right)
= $ \ans{the angle between the two vectors}

\vspace{0.12in}
\textbf{Perpendicular vectors.} Two nonzero vectors are perpendicular if and only if their
dot product $= $ \ans{$0$} \quad (because $\cos 90^\circ = $ \ans{$0$}).

\vspace{0.12in}
\textbf{Example.} $\vec{u}=\langle 1,2\rangle$, $\vec{w}=\langle 4,-2\rangle$.

\vspace{0.06in}
\hspace{1em} $\vec{u}\cdot\vec{w} = (1)(4) + (2)(-2) = $ \ans{$0$}

\hspace{1em} Are $\vec{u}$ and $\vec{w}$ perpendicular? \ans{Yes, because the dot product equals 0.}

\vspace{0.1in}
\textbf{Law of Cosines application.} If $\vec{r} = \vec{u}+\vec{w}$, the triangle with sides
$\|\vec{u}\|$, $\|\vec{w}\|$, $\|\vec{r}\|$ satisfies:
$\|\vec{r}\|^2 = \|\vec{u}\|^2 + \|\vec{w}\|^2 - 2\|\vec{u}\|\,\|\vec{w}\|\cos\theta$

where $\theta$ is the angle \ans{between $\vec{u}$ and $\vec{w}$}.
\end{notesbox}

\end{document}
